\documentclass[11pt,a4paper]{article}

\usepackage[T1]{fontenc}
\usepackage[spanish]{babel}
\usepackage[margin=2.5cm]{geometry}
\usepackage{graphicx}
\usepackage{booktabs}
\usepackage{tabularx}
\usepackage[hidelinks]{hyperref}

\newcolumntype{Y}{>{\raggedright\arraybackslash}X}
\newcommand{\pendiente}[1]{\textit{[Pendiente: #1]}}

\title{TA-PFC-E1 · Proyecto Fin de Curso\\
Visión, alcance, partes interesadas, contexto y plan de elicitación\\
Sistema de Gestión de Microcréditos (SGMC)}
\author{Universidad Técnica Estatal de Quevedo\\
Facultad de Ciencias de la Computación · Carrera de Software\\
Ingeniería de Requerimientos (ISR-401)\\[0.6em]
CHAVARRIA CUENCA TAHINY MEL\\
RIVERA SUAREZ JONATHAN JAVIER\\
ROCHINA ROCHINA FREDDY DAVID}
\date{SPA 2026--2027}

\begin{document}
\maketitle

\section{Organización y problema}
\textbf{Organización:} \pendiente{nombre oficial y ubicación en Quevedo o Los Ríos}.

\textbf{Problema acordado:} \pendiente{proceso actual, dificultad observada, impacto y acuerdo con la organización sobre el problema abordable con software}.

\textbf{Respaldo:} \pendiente{identificadores de la carta de aceptación firmada y del registro del primer contacto en \path{evidencias/MANIFIESTO.md}}.

\section{Visión y alcance}
La delimitación del SGMC se enmarca en el proceso de ingeniería de requisitos de la norma ISO/IEC/IEEE 29148 \cite{iso29148}. La visión, los objetivos y el alcance se presentan como propuesta para validar con la organización.

\subsection{Declaración de visión}
El Sistema de Gestión de Microcréditos (SGMC) permitirá al personal que recibe, evalúa y da seguimiento a las solicitudes consultar en un solo expediente el estado de cada trámite, sus documentos, las decisiones adoptadas y los pagos registrados. Así, la organización podrá identificar qué falta para avanzar, comunicar el estado del trámite a la persona solicitante y reconstruir el historial de cada crédito desde la recepción de la solicitud hasta el cierre de la obligación.

\subsection{Objetivos de negocio}
Las metas se proponen para el piloto. La organización deberá validar su viabilidad y el equipo deberá levantar la línea base antes de medir el cambio.

\begin{enumerate}
  \item \textbf{Reducir el tiempo de tramitación.} Reducir en un 20\,\% la mediana de días hábiles entre la recepción de una solicitud completa y la comunicación de su decisión al finalizar el tercer mes del piloto. Comparar las fechas de recepción y comunicación con las del período de línea base.
  \item \textbf{Reducir las devoluciones por información incompleta.} Reducir en un 25\,\% la proporción de solicitudes ingresadas que se devuelven al menos una vez por datos o documentos faltantes durante el tercer mes del piloto, respecto a la línea base. Verificar el cambio con el motivo registrado para cada devolución.
  \item \textbf{Mejorar la trazabilidad de los pagos.} Lograr que al menos el 95\,\% de los pagos registrados durante el tercer mes del piloto estén asociados al crédito correspondiente y cuenten con fecha, monto y referencia a un comprobante. Verificarlo mediante una revisión de los registros de pago del período.
\end{enumerate}

\subsection{Alcance y restricciones}
\textbf{Incluye.} Registrar solicitudes y sus documentos, revisar si la información está completa, registrar devoluciones y consultar el estado del trámite. Durante la evaluación, el personal autorizado documentará la decisión y las condiciones aprobadas. Para los créditos formalizados, se registrará la confirmación del desembolso, el calendario de cuotas validado, los pagos informados y el cierre de la obligación. El historial identificará la fecha, el rol responsable y el motivo de los cambios de estado.

\textbf{Excluye.} La aprobación automática, la ejecución de desembolsos o cobros, la contabilidad general y la integración automática con bancos o burós de crédito quedan fuera del alcance propuesto. El alcance contempla el registro de operaciones realizadas por los canales autorizados de la organización.

\textbf{Restricciones.} La decisión crediticia corresponde al personal autorizado. Las reglas de evaluación, los estados del trámite y el cálculo o aplicación de cuotas y pagos requieren confirmación de la organización antes de su incorporación al sistema. El acceso a expedientes y comprobantes se limitará a los roles autorizados. El repositorio público solo incluirá evidencias anonimizadas y el manifiesto.

\section{Partes interesadas}
La tabla~\ref{tab:partes} registra las partes interesadas por rol y su estrategia de participación.

\begin{table}[htbp]
\centering
\caption{Registro de partes interesadas}
\label{tab:partes}
\small
\begin{tabularx}{\linewidth}{@{}lYYYYY@{}}
\toprule
ID & Rol & Interés & Influencia & Disponibilidad & Estrategia \\
\midrule
P1 & \pendiente{rol} & --- & --- & --- & --- \\
P2 & \pendiente{rol} & --- & --- & --- & --- \\
P3 & \pendiente{rol} & --- & --- & --- & --- \\
P4 & \pendiente{rol} & --- & --- & --- & --- \\
P5 & \pendiente{rol} & --- & --- & --- & --- \\
P6 & \pendiente{rol} & --- & --- & --- & --- \\
P7 & \pendiente{rol} & --- & --- & --- & --- \\
P8 & \pendiente{rol} & --- & --- & --- & --- \\
\bottomrule
\end{tabularx}
\end{table}

\subsection{Matriz poder--interés}
\IfFileExists{figuras/matriz-poder-interes.pdf}{%
  La figura~\ref{fig:matriz} presenta la ubicación de los roles por poder e interés.
  \begin{figure}[htbp]
    \centering
    \includegraphics[width=0.85\linewidth,height=0.65\textheight,keepaspectratio]{figuras/matriz-poder-interes.pdf}
    \caption{Matriz poder--interés}
    \label{fig:matriz}
  \end{figure}
}{\pendiente{incorporar la matriz exportada y guardar su fuente editable en \texttt{modelos/}}.}

\pendiente{justificar la ubicación de los roles y su relación con las estrategias de participación}.

\section{Contexto}

\subsection{Diagrama de contexto}
\IfFileExists{figuras/diagrama-contexto.pdf}{%
  La figura~\ref{fig:contexto} muestra el límite del sistema y sus intercambios con las entidades externas.
  \begin{figure}[htbp]
    \centering
    \includegraphics[width=0.9\linewidth,height=0.65\textheight,keepaspectratio]{figuras/diagrama-contexto.pdf}
    \caption{Diagrama de contexto del SGMC}
    \label{fig:contexto}
  \end{figure}
}{\pendiente{incorporar el diagrama exportado y guardar su fuente editable en \texttt{modelos/}}.}

\pendiente{explicar el límite acordado, las entidades externas y los intercambios identificados}.

\subsection{Ocho perspectivas del sílabo}
La tabla~\ref{tab:perspectivas} vincula cada perspectiva con un hallazgo de la organización y su respaldo.

\begin{table}[htbp]
\centering
\caption{Perspectivas del sílabo y hallazgos}
\label{tab:perspectivas}
\begin{tabularx}{\linewidth}{@{}lYYY@{}}
\toprule
N.º & Perspectiva del sílabo & Hallazgo & Fuente o evidencia \\
\midrule
1 & \pendiente{nombre} & --- & --- \\
2 & \pendiente{nombre} & --- & --- \\
3 & \pendiente{nombre} & --- & --- \\
4 & \pendiente{nombre} & --- & --- \\
5 & \pendiente{nombre} & --- & --- \\
6 & \pendiente{nombre} & --- & --- \\
7 & \pendiente{nombre} & --- & --- \\
8 & \pendiente{nombre} & --- & --- \\
\bottomrule
\end{tabularx}
\end{table}

\section{Plan de elicitación}

\subsection{Técnicas y justificación por rol}
La elicitación parte de las fuentes de requisitos y del conocimiento que se quiere recuperar~\cite{pohl}. Se propone usar entrevistas para explorar decisiones y excepciones, un cuestionario para contrastar la frecuencia percibida de dificultades y observación para documentar la secuencia del proceso. La asignación a roles requiere validación con la organización.
La tabla~\ref{tab:tecnicas} relaciona los roles con la información que se busca, la técnica seleccionada y el instrumento correspondiente.

\begin{table}[htbp]
\centering
\caption{Selección de técnicas de elicitación}
\label{tab:tecnicas}
\small
\noindent\begin{tabularx}{\linewidth}{@{}lYYYY@{}}
\toprule
Rol & Información buscada & Técnica & Justificación por fuente y tipo de conocimiento & Instrumento \\
\midrule
P1 & --- & --- & --- & --- \\
P2 & --- & --- & --- & --- \\
P3 & --- & --- & --- & --- \\
P4 & --- & --- & --- & --- \\
P5 & --- & --- & --- & --- \\
P6 & --- & --- & --- & --- \\
P7 & --- & --- & --- & --- \\
P8 & --- & --- & --- & --- \\
\bottomrule
\end{tabularx}
\end{table}

\subsection{Cronograma de sesiones}
\textit{[Pendiente: registrar una fila por sesión planificada en la tabla~\ref{tab:sesiones}]}.

\begin{table}[htbp]
\centering
\caption{Cronograma de sesiones de elicitación}
\label{tab:sesiones}
\begin{tabularx}{\linewidth}{@{}YYYYY@{}}
\toprule
Fecha y duración & Roles convocados & Técnica e instrumento & Responsable & Objetivo de la sesión \\
\midrule
\pendiente{fecha} & --- & --- & --- & --- \\
\bottomrule
\end{tabularx}
\end{table}

\subsection{Instrumentos}
\begin{itemize}
  \item Guion de entrevista: \pendiente{confirmar que el documento de preguntas existente constituye el guion completo e indicar su ruta definitiva}.
  \item Cuestionario: \path{instrumentos/Cuestionario.md} (borrador por validar).
  \item Guía de observación o de taller: \path{instrumentos/Guia_observacion.md} (borrador por validar).
\end{itemize}

\section{Ética y tratamiento de datos}

\subsection{Modelo de consentimiento informado}
\textbf{Propósito.} La sesión académica busca comprender el proceso de microcrédito para definir el alcance del SGMC. Se recogerán notas sobre las actividades y decisiones del rol.

\textbf{Participación.} La participación es voluntaria. Se puede omitir preguntas o finalizar la sesión. El informe identificará las respuestas por código de rol. Se solicitará autorización específica para grabar.

\textbf{Custodia.} Los registros identificables se guardarán en la carpeta institucional privada. \pendiente{completar y comunicar el responsable, el canal de contacto y el plazo de conservación antes de aplicar el formulario}.

\textbf{Acepta participar:} Sí / No. \quad
\textbf{Autoriza grabación:} Sí / No / No aplica.

\textbf{Fecha y firma:} \rule{5cm}{0.4pt}.

\subsection{Anonimización y custodia}
El tratamiento de datos se enmarca en la legislación ecuatoriana y su reglamento~\cite{lopdp,reg904}. Para gestionar riesgos sobre la confidencialidad e integridad de las evidencias se toma como referencia ISO/IEC 27001. Para delimitar el acceso a los registros privados se considera la orientación sobre controles de ISO/IEC 27002~\cite{iso27001,iso27002}.
\pendiente{detallar cómo el equipo recopilará, anonimizará y revisará la información, y quién administrará el acceso a la carpeta institucional privada}.

\par La carta firmada y las evidencias con datos personales se custodian en la carpeta institucional privada compartida con el docente. El repositorio incluye las evidencias anonimizadas y el manifiesto. Cada registro indica nombre del archivo, fecha, duración y huella SHA-256.

\section{Gestión del proyecto}
La planificación del equipo, la comunicación con las partes interesadas y el seguimiento de los riesgos se organizan con apoyo del marco de gestión del PMBOK~\cite{pmbok}.

\subsection{Roles del equipo}


\noindent\textbf{CHAVARRIA CUENCA TAHINY MEL.} Coordinación y enlace con la organización. Coordina el acuerdo sobre el problema, la visión y el alcance. Da seguimiento al cronograma hasta la semana 17, a los acuerdos de comunicación y al registro de riesgos.

\noindent\textbf{RIVERA SUAREZ JONATHAN JAVIER.} Coordina la elicitación. Prepara el cronograma de sesiones y revisa los tres instrumentos. Gestiona el consentimiento informado, la anonimización y el manifiesto de evidencias.

\noindent\textbf{ROCHINA ROCHINA FREDDY DAVID.} Análisis de partes interesadas y modelado. Mantiene el registro de roles y sus estrategias de participación. Prepara la matriz poder--interés, el diagrama de contexto y la tabla de las ocho perspectivas.


\subsection{Cronograma hasta la semana 17}
\textit{[Pendiente: completar la tabla~\ref{tab:proyecto} con las actividades, los hitos y sus responsables hasta la semana 17]}.

\begin{table}[htbp]
\centering
\caption{Cronograma del proyecto}
\label{tab:proyecto}
\begin{tabularx}{\linewidth}{@{}lYYY@{}}
\toprule
Semana & Actividad & Responsable & Hito o producto \\
\midrule
\pendiente{semana} & --- & --- & --- \\
\bottomrule
\end{tabularx}
\end{table}

\subsection{Comunicación con la organización}
\pendiente{definir los roles de contacto, el canal, la frecuencia, los responsables y el registro de acuerdos}.

\subsection{Registro de riesgos}
\textit{[Pendiente: completar la tabla~\ref{tab:riesgos} con riesgos específicos de la organización y del proyecto]}.

\begin{table}[htbp]
\centering
\caption{Registro de riesgos del proyecto}
\label{tab:riesgos}
\small
\begin{tabularx}{\linewidth}{@{}lYYYYY@{}}
\toprule
ID & Riesgo y causa & Probabilidad & Impacto & Respuesta & Responsable \\
\midrule
R1 & \pendiente{riesgo} & --- & --- & --- & --- \\
\bottomrule
\end{tabularx}
\end{table}

\section{Declaración de uso de inteligencia artificial}
Se utilizó ChatGPT mediante Codex como apoyo para preparar la plantilla LaTeX, los borradores de instrumentos y la integración de las referencias. \pendiente{completar la versión consultable, registrar otros usos si los hubo y describir cómo el equipo verificó lo obtenido}.

\bibliographystyle{IEEEtran}
\bibliography{referencias}

\end{document}
